\documentclass[11pt,oneside,openany]{book}
\usepackage[T1]{fontenc}
\usepackage[utf8]{inputenc}
\usepackage[sc]{mathpazo}
\renewcommand{\sfdefault}{lmss}
\renewcommand{\ttdefault}{lmtt}
\usepackage{amsmath,amssymb,amsthm,mathtools}
\usepackage[paperwidth=7.5in,paperheight=10in,inner=.85in,outer=.85in,top=.85in,bottom=.85in,headheight=15pt,headsep=20pt]{geometry}
\usepackage{microtype}
\usepackage{xcolor,graphicx,booktabs,array,longtable}
\newcolumntype{L}[1]{>{\raggedright\arraybackslash}p{#1}}
\usepackage{tikz,pgfplots}
\usetikzlibrary{arrows.meta,calc,patterns,decorations.markings}
\pgfplotsset{compat=1.18}
\definecolor{ink}{HTML}{17334B}
\definecolor{teal}{HTML}{137C83}
\definecolor{copper}{HTML}{AE5932}
\definecolor{soft}{HTML}{EEF3F5}
\pgfplotsset{every axis/.append style={axis line style={black!65},tick style={black!60},tick label style={font=\small},label style={font=\small},legend style={font=\small,draw=none,fill=white},grid style={black!10},major grid style={black!10},line width=.6pt},every axis plot/.append style={line width=1.15pt}}
\usepackage{enumitem,needspace,etoolbox}
\setlist[enumerate]{topsep=5pt,itemsep=5pt,parsep=0pt}
\setlist[itemize]{topsep=4pt,itemsep=3pt,parsep=0pt}
\usepackage{titlesec,fancyhdr}
\titleformat{\chapter}[display]{\normalfont\color{ink}}{\large\sffamily\MakeUppercase{Chapter \thechapter}}{12pt}{\huge\bfseries\raggedright\hyphenpenalty=10000\exhyphenpenalty=10000}[\vspace{12pt}\titlerule]
\titlespacing*{\chapter}{0pt}{-10pt}{24pt}
\titleformat{\section}{\large\bfseries\color{ink}}{\thesection}{.65em}{}
\titleformat{\subsection}{\normalsize\bfseries\color{ink}}{\thesubsection}{.6em}{}
\titlespacing*{\section}{0pt}{18pt plus 3pt minus 2pt}{7pt}
\titlespacing*{\subsection}{0pt}{12pt plus 2pt minus 2pt}{5pt}
\pagestyle{fancy}
\fancyhf{}
\fancyhead[L]{\small\nouppercase{\leftmark}}
\fancyfoot[C]{\small\thepage}
\renewcommand{\headrulewidth}{.3pt}
\renewcommand{\chaptermark}[1]{\markboth{\thechapter\quad #1}{}}
\fancypagestyle{plain}{\fancyhf{}\fancyfoot[C]{\small\thepage}\renewcommand{\headrulewidth}{0pt}}
\usepackage[font=small,labelfont=bf,labelsep=period]{caption}
\usepackage{float,placeins}
\floatplacement{figure}{htbp}
\setcounter{topnumber}{3}
\setcounter{bottomnumber}{2}
\renewcommand{\topfraction}{.9}
\renewcommand{\textfraction}{.08}
\renewcommand{\floatpagefraction}{.8}
\usepackage{makeidx}
\makeindex
\usepackage[pdfencoding=auto,psdextra,colorlinks=true,linkcolor=ink,citecolor=teal,urlcolor=teal,bookmarksnumbered=true]{hyperref}
\usepackage{bookmark}
\hypersetup{pdftitle={Discrete Information Geometry},pdfsubject={Finite differences, probability, and information flows}}
\urlstyle{same}
\newtheorem{theorem}{Theorem}[chapter]
\newtheorem{proposition}[theorem]{Proposition}
\newtheorem{lemma}[theorem]{Lemma}
\newtheorem{corollary}[theorem]{Corollary}
\theoremstyle{definition}
\newtheorem{definition}[theorem]{Definition}
\newtheorem{example}{Example}[chapter]
\BeforeBeginEnvironment{example}{\par\addvspace{4pt}\noindent\begin{minipage}{\linewidth}}
\AfterEndEnvironment{example}{\end{minipage}\par\addvspace{4pt}}
\theoremstyle{remark}
\newtheorem*{remark}{Remark}
\newenvironment{solution}{\par\smallskip\noindent\textit{Solution.}\ }{\par\medskip}
\newenvironment{exercises}{\begin{enumerate}[label=\thechapter.\arabic*,ref=\thechapter.\arabic*,leftmargin=2.8em,itemsep=6pt]}{\end{enumerate}}
\newenvironment{answers}{\begin{description}[leftmargin=2.9em,labelwidth=2.5em,labelsep=.4em,itemsep=8pt,font=\normalfont\bfseries]}{\end{description}}
\newcommand{\R}{\mathbb{R}}
\newcommand{\E}{\mathbb{E}}
\newcommand{\Var}{\operatorname{Var}}
\newcommand{\Cov}{\operatorname{Cov}}
\newcommand{\KL}[2]{D\!\left(#1\,\middle\|\,#2\right)}
\newcommand{\grad}{\operatorname{grad}}
\newcommand{\arcosh}{\operatorname{arcosh}}
\newcommand{\sech}{\operatorname{sech}}
\newcommand{\softmax}{\operatorname{softmax}}
\newcommand{\diag}{\operatorname{diag}}
\newcommand{\dd}{\,\mathrm{d}}
\newcommand{\ip}[2]{\left\langle #1,#2\right\rangle}
\newcommand{\p}{\mathbf{p}}
\newcommand{\q}{\mathbf{q}}
\newcommand{\f}{\mathbf{f}}
\newcommand{\uvec}{\mathbf{u}}
\newcommand{\vvec}{\mathbf{v}}
\newcommand{\ones}{\mathbf{1}}
\newcommand{\simp}{\Delta^{n-1}}
\newcommand{\simpi}{\Delta^{n-1}_{\!\circ}}
\numberwithin{equation}{chapter}
\setlength{\parindent}{1.2em}
\setlength{\parskip}{0pt}
\linespread{1.07}
\setlength{\emergencystretch}{2em}
\clubpenalty=10000
\widowpenalty=10000
\displaywidowpenalty=10000
\raggedbottom
\setcounter{tocdepth}{1}
\makeatletter
\renewcommand*\l@section{\@dottedtocline{1}{1.5em}{3.2em}}
\makeatother
\setcounter{secnumdepth}{2}
\providecommand{\tightlist}{\setlength{\itemsep}{0pt}\setlength{\parskip}{0pt}}
\providecommand{\passthrough}[1]{#1}
\newcommand{\labheading}{\FloatBarrier\section*{Laboratory investigation}\addcontentsline{toc}{section}{Laboratory investigation}}

\begin{document}
\begin{titlepage}\color{ink}\centering\vspace*{1.0in}{\small\sffamily PROBABILITY AND MATHEMATICAL MODELLING\par}\vspace{.45in}{\Huge\bfseries Discrete Information\\[.12in]Geometry\par}\vspace{.35in}{\Large Finite Differences, Probability,\\and Information Flows\par}\vspace{1cm}{\Large Chapter 10\par}\vfill\begin{tikzpicture}[scale=1.05]\draw[teal,line width=1pt](0,0)--(4,0)--(2,3.464)--cycle;\foreach \a in {0,...,6}{\foreach \b in {0,...,6}{\pgfmathtruncatemacro{\c}{6-\a-\b}\ifnum\c>-1\fill[teal] ({(\b+\c/2)*4/6},{\c*3.464/6}) circle(1.5pt);\fi}}\end{tikzpicture}\vfill{\large A course in ten chapters\par}\vspace{.3in}{\small\url{https://bobovski66.github.io/IntroDiscreteInfoGeom/}\par}\end{titlepage}
\setcounter{chapter}{9}

\chapter{Graph divergence, diffusion, and information flow}
A network can transport something, record something, or do both. A random walker changes location on a graph. A sensor converts a detailed physical state into an imperfect report. An observer may see only a group label rather than the original category. Each operation transforms a probability distribution.

The tools developed throughout the course now meet in one setting. Edge differences produce currents; divergence gives conservation; entropy measures uncertainty; relative entropy and spherical distance measure distinctions between models. We will examine which quantities a network preserves, which it changes, and what becomes impossible to recover after observations are grouped.

\section{From an edge gradient to diffusion}\label{ch10-from-an-edge-gradient-to-diffusion}

Let a finite undirected graph have vertices \(i\) and symmetric nonnegative edge weights \(c_{ij}=c_{ji}\). A positive weight means that the edge permits transport. Let \(p_i\) be the probability stored at vertex \(i\).

The graph gradient on an oriented edge is \[(\nabla_G p)_{ij}=p_j-p_i.\] It is positive when the probability is higher at the destination. Diffusion moves down this difference, from larger concentration to smaller concentration. With step parameter \(\tau\geq0\), define the outward current \[J_{ij}=\tau c_{ij}(p_i-p_j)=-\tau c_{ij}(\nabla_Gp)_{ij}.\] The current reverses sign when the edge orientation is reversed.

The \textbf{graph Laplacian} is the vertex operator \[\boxed{(Lp)_i=\sum_{j\sim i}c_{ij}(p_i-p_j).}\] It adds the weighted differences between the value at a vertex and the values at its neighbors. With our outward-divergence convention, \(\operatorname{div}J=\tau Lp\). The continuity equation becomes \[\boxed{\Delta p=-\tau Lp,\qquad p^+=(I-\tau L)p.}\]

This is a complete discrete diffusion equation. The time difference on the left records one update; the graph operator on the right measures imbalance across adjacent locations.

\section{A step-size condition preserves probabilities}\label{ch10-a-step-size-condition-preserves-probabilities}

Let \(d_i=\sum_{j\sim i}c_{ij}\) be the weighted degree of vertex \(i\). Expanding the update gives \[p_i^+=(1-\tau d_i)p_i+\tau\sum_{j\sim i}c_{ij}p_j.\] A sufficient condition for all coefficients to be nonnegative is \[0\leq\tau d_i\leq1\quad\text{for every vertex }i.\] When the graph has a positive edge, this is ensured by \[\tau\leq\frac1{\max_i d_i}.\] On a graph with no edges, \(L=0\) and nothing changes.

Symmetry of the edge weights makes \(K=I-\tau L\) symmetric. Its rows and columns each sum to one. Under the step-size condition it is therefore \textbf{doubly stochastic}: both a valid transition matrix and a table whose row sums are one.

The algebra explains why excessive steps are problematic. If \(\tau d_i>1\), the retained-mass coefficient at vertex \(i\) is negative. Starting with all mass at that vertex can then produce a negative probability. A time step belongs to the model, just as correction strength did in Chapter 2.

\section{The three-node path gives an exact example}\label{ch10-the-three-node-path-gives-an-exact-example}

For the path \(A\)--\(B\)--\(C\) with unit edge weights, \[L=\begin{pmatrix}1&-1&0\\-1&2&-1\\0&-1&1\end{pmatrix}.\] The largest degree is two, so \(0\leq\tau\leq1/2\) preserves nonnegativity. Taking \(\tau=1/4\) recovers the matrix from Chapter 4: \[K=\begin{pmatrix}3/4&1/4&0\\1/4&1/2&1/4\\0&1/4&3/4\end{pmatrix}.\]

Starting from \((1,0,0)\) gives first \((3/4,1/4,0)\) and then \((5/8,5/16,1/16)\). The uniform vector \(u=(1/3,1/3,1/3)\) is stationary.

We can solve every trajectory by splitting deviations into two patterns: \[a(1,0,-1)+b(1,-2,1).\] Both patterns have coordinate sum zero. Every probability vector can be written \[p_0=u+a(1,0,-1)+b(1,-2,1),\] where \[a=\frac{p_{A,0}-p_{C,0}}2,\qquad
b=\frac{p_{A,0}+p_{C,0}-2p_{B,0}}6.\]

Direct matrix multiplication shows \[K(1,0,-1)=\frac34(1,0,-1),\qquad
K(1,-2,1)=\frac14(1,-2,1).\] The first imbalance shrinks by three quarters each step; the second shrinks by one quarter. Linearity then gives \[\boxed{
p_n=u+a(3/4)^n(1,0,-1)+b(1/4)^n(1,-2,1).
}\] Both powers tend to zero, proving convergence to uniform from every initial distribution.

For \(p_0=(1,0,0)\), the coefficients are \(a=1/2\) and \(b=1/6\). The slower pattern measures imbalance between the two ends of the path. This decomposition shows how a network can have several rates of relaxation without requiring us to solve a new recurrence at each vertex.

\begin{example}[A centered initial state relaxes at the faster rate]

Start the three-node walk at \(p_0=(0,1,0)\). Its end imbalance coefficient is \(a=0\), while \(b=-1/3\). Therefore \[p_n=\left(\frac13-\frac13(1/4)^n,\;
\frac13+\frac23(1/4)^n,\;
\frac13-\frac13(1/4)^n\right).\] The endpoint probabilities remain equal by symmetry. Only the faster pattern is present. Starting at one endpoint excites both patterns and retains the slower factor \((3/4)^n\).

\end{example}

A network's relaxation rate depends both on its transition rule and on which patterns the initial state contains. The exact formula makes that dependence visible. A plot of one initial state may miss a slowly decaying pattern that another initial state excites.

For this example, the error at the middle vertex is \((2/3)(1/4)^n\). Requiring it to be below \(0.01\) reduces to a comparison of powers, just like the scalar stopping questions in Chapter 2. The network problem has been decomposed into scalar difference equations.

These patterns are often called modes. A mode is a vector pattern reproduced up to a multiplier by the update. No new notation is needed to verify one: multiply the vector by the matrix and compare its coordinates with the proposed multiple.

\section{Entropy rises under a doubly stochastic update}\label{ch10-entropy-rises-under-a-doubly-stochastic-update}

Let \(h(x)=-x\ln x\), with \(h(0)=0\). Chapter 5 proved its concavity. If \(K\) is doubly stochastic, each output coordinate is a weighted average of input coordinates because its row sums to one: \[h\left(\sum_jK_{ij}p_j\right)\geq\sum_jK_{ij}h(p_j).\] Summing over output coordinates and using the column sums gives \[\boxed{H(Kp)\geq H(p).}\] Thus the nonnegative entropy increment \[\sigma_n=H(p_{n+1})-H(p_n)\] is guaranteed for this diffusion rule.

For the three-node example, \(H(p_0)=0\), \(H(p_1)\approx0.562335\), and the exact trajectory approaches entropy \(\ln3\). Telescoping yields \[\sum_{n=0}^{T-1}\sigma_n=H(p_T)-H(p_0)\leq\ln3-H(p_0).\] The total possible entropy gain is bounded.

Doubly stochastic does not by itself imply convergence. A permutation matrix is doubly stochastic and preserves entropy exactly. The two-vertex diffusion rule at \(\tau=1\) simply swaps its coordinates. A convergence claim needs information about the dynamics beyond entropy monotonicity.

If the graph is disconnected, mass cannot travel between components. Its stationary behavior is constrained by the mass initially present in each component. Global uniformity is not generally reachable from every initial state.

\section{A general observation channel has another interpretation}\label{ch10-a-general-observation-channel-has-another-interpretation}

A \textbf{channel} is a column-stochastic matrix \(K\) whose entry \(K_{yi}\) is the probability of reporting output \(y\) given input label \(i\). The output and input label sets may have different sizes. An input distribution \(p\) produces output distribution \[P_y=\sum_iK_{yi}p_i.\] This is the same matrix arithmetic as transport, now describing observation.

A deterministic grouping is a special channel. If labels one and two are reported as group \(A\), while labels three and four are reported as group \(B\), then \[K=\begin{pmatrix}1&1&0&0\\0&0&1&1\end{pmatrix}.\] The output retains group totals and discards the original label within each group.

An arbitrary channel need not increase entropy. The two-input channel that always reports the same single symbol sends every input to a certain output of entropy zero. Even a square channel can do this, for example \[K=\begin{pmatrix}1&1\\0&0\end{pmatrix}.\] It is column-stochastic but not doubly stochastic.

There is, however, a general statement about distinguishing two input models: passing both through the same channel cannot increase their KL divergence or their Fisher--Rao distance. The reference to the \textbf{same} channel is essential.

The observational interpretation differs from the physical transport interpretation in what is claimed to happen to an individual state. A grouping report does not require the original categories to merge physically. It records only a coarser label. The fine state can remain as detailed as before while the observer loses access to some of its distinctions.

This distinction matters when reading data-processing inequalities. They compare the probability tables available before and after the chosen recording operation. They do not claim that every internal feature of the system has ceased to exist. A more informative measurement scheme may retain distinctions that another scheme erases.

Observation rules can be composed. If \(K\) converts a fine label to an intermediate report and \(R\) converts that report to a final report, the combined channel is \(RK\). Its entries are nonnegative, and its column sums are one because \[\sum_z(RK)_{zi}
=\sum_yK_{yi}\sum_zR_{zy}
=\sum_yK_{yi}=1.\] Two stages of processing therefore obey the same framework as one. Data processing can be applied successively to compare the distinctions visible at each stage.

\section{The log-sum inequality organizes the proof}\label{ch10-the-log-sum-inequality-organizes-the-proof}

For nonnegative lists \(a_i,b_i\), let \(A=\sum_i a_i\) and \(B=\sum_i b_i\). When the relevant sums are positive and finite, define normalized tables \(\alpha_i=a_i/A\) and \(\beta_i=b_i/B\). Expanding gives \[\sum_i a_i\ln\frac{a_i}{b_i}
=A\ln\frac AB+A D(\alpha\|\beta).\] Gibbs' inequality therefore implies \[\boxed{A\ln(A/B)\leq\sum_i a_i\ln(a_i/b_i).}\] This is the \textbf{log-sum inequality}. It says grouping terms can reduce a directed logarithmic discrepancy. The established zero and infinity conventions extend it to boundary cases.

Apply the inequality separately to each output label, with \[a_i=K_{yi}p_i,\qquad b_i=K_{yi}q_i.\] Their sums are \((Kp)_y\) and \((Kq)_y\). For nonzero channel entries, the factor \(K_{yi}\) cancels inside the ratio. Summing over outputs gives \[\begin{aligned}
D(Kp\|Kq)
&\leq\sum_{y,i}K_{yi}p_i\ln\frac{p_i}{q_i}\\
&=\sum_i p_i\ln\frac{p_i}{q_i}\sum_yK_{yi}\\
&=D(p\|q).
\end{aligned}\] This is the \textbf{data-processing inequality}.

The proof uses only finite sums and Gibbs' inequality. Its meaning is observational: a common reporting mechanism cannot make two models more distinguishable by this measure. It can preserve a distinction or erase some of it.

\section{Relative entropy to stationarity cannot increase}\label{ch10-relative-entropy-to-stationarity-cannot-increase}

If \(\pi\) is a stationary distribution, then \(K\pi=\pi\). Data processing gives \[D(Kp\|\pi)=D(Kp\|K\pi)\leq D(p\|\pi).\] Thus \(D(p_n\|\pi)\) is a nonincreasing potential whenever its initial value is finite.

This generalizes the Lyapunov viewpoint from Chapter 6 to Markov transport. It does not automatically provide the geometric contraction factor proved for the mixture rule. Some channels preserve the divergence. A separate estimate is needed to obtain a numerical convergence rate.

When \(\pi\) is uniform, the relation \(D(p\|\pi)=\ln k-H(p)\) turns divergence decrease into entropy increase. With a nonuniform stationary distribution, relative entropy can decrease without ordinary entropy increasing. The correct target is determined by the transition mechanism.

An asymmetric two-location channel shows why stationarity, rather than uniformity, is the appropriate reference in general. Take \[K=\begin{pmatrix}3/4&1/2\\1/4&1/2\end{pmatrix}.\] Its stationary distribution is \(\pi=(2/3,1/3)\), from the two-location calculation in Chapter 4. Starting at uniform gives the next state \((5/8,3/8)\). Entropy decreases because the new table is less balanced. Nevertheless, \[D((5/8,3/8)\|\pi)\leq D((1/2,1/2)\|\pi).\] The process is approaching its own stationary state rather than maximizing uncertainty over all possible states.

A model can also have several stationary distributions. On a graph with disconnected components, any compatible allocation of total mass among component equilibria remains stationary. A target must then account for those conserved component masses. Measuring divergence to an unreachable reference could misrepresent what the dynamics can accomplish.

In the connected symmetric three-node example, the exact solution identifies a unique stationary distribution and proves convergence. In a new graph, begin by checking connectivity, retained mass, and possible periodicity before generalizing that conclusion.

\section{Square-root distance also contracts}\label{ch10-square-root-distance-also-contracts}

For each output label, Cauchy--Schwarz gives \[\sqrt{\left(\sum_iK_{yi}p_i\right)
\left(\sum_iK_{yi}q_i\right)}
\geq\sum_iK_{yi}\sqrt{p_iq_i}.\] Add over outputs and use the channel column sums: \[B(Kp,Kq)\geq B(p,q).\] The output distributions have at least as much overlap. Since inverse cosine decreases as its argument increases, \[\boxed{d_{\mathrm{FR}}(Kp,Kq)\leq d_{\mathrm{FR}}(p,q).}\] The Hellinger distance also cannot increase because its square is \(1-B\).

This result links the sphere geometry to observation. A channel is not merely a linear map between numerical lists: it has a definite effect on the statistical distances we constructed.

\section{Grouping shows exactly where a distinction is lost}\label{ch10-grouping-shows-exactly-where-a-distinction-is-lost}

Take \[p=(0.1,0.4,0.1,0.4),\qquad
q=(0.4,0.1,0.4,0.1).\] The input models differ within both pairs of labels. Their divergence is \[D(p\|q)=0.6\ln4>0.\] After grouping \(\{1,2\}\) and \(\{3,4\}\), both output tables are \((1/2,1/2)\). The output divergence and spherical distance are zero.

\begin{figure}[htbp]\centering
\resizebox{.98\linewidth}{!}{\begin{tikzpicture}[every node/.style={font=\small},>=Stealth]
\node[anchor=west] at (0,4.5){Input label};\node[anchor=west] at (2,4.5){$q_i$};\node[anchor=west] at (3.1,4.5){$p_i$};
\node[anchor=west] at (7,4.5){Recorded group};
\foreach \i/\y/\q/\p in {1/3.5/.4/.1,2/2.5/.1/.4,3/1.5/.4/.1,4/.5/.1/.4}{
\node at (.5,\y){\i};\node at (2.3,\y){\q};\node at (3.4,\y){\p};}
\draw[->,teal,thick] (4,3.5)--(7,3);\draw[->,teal,thick] (4,2.5)--(7,3);
\draw[->,copper,thick] (4,1.5)--(7,1);\draw[->,copper,thick] (4,.5)--(7,1);
\node[anchor=west,align=left] at (7.2,3){Group $\{1,2\}$\\Both models: $1/2$};
\node[anchor=west,align=left] at (7.2,1){Group $\{3,4\}$\\Both models: $1/2$};
\node[anchor=west] at (0,-.5){$D(q\Vert p)=0.6\ln4\quad\longrightarrow\quad D(Kq\Vert Kp)=0$};
\end{tikzpicture}}
\caption{Grouping can erase a distinction carried entirely by the probabilities inside each group. Both recorded output tables are identical.}\end{figure}

We can account for the lost divergence exactly. For each group \(g\), let \[P_g=\sum_{i\in g}p_i,\qquad Q_g=\sum_{i\in g}q_i.\] When the group masses are positive, define conditional distributions within that group: \[p(i\mid g)=p_i/P_g,\qquad q(i\mid g)=q_i/Q_g.\] The conditional probability describes the fine label after the group is known.

Factor the ratio as \[\frac{p_i}{q_i}=\frac{P_g}{Q_g}
\frac{p(i\mid g)}{q(i\mid g)}.\] Taking logarithms, multiplying by \(p_i\), and summing first within groups gives \[\boxed{
D(p\|q)=D(P\|Q)+\sum_g P_gD(p(\cdot\mid g)\|q(\cdot\mid g)).
}\] For positive tables this is an exact identity. Terms of zero group mass are omitted under the usual conventions.

The first term is the distinction still visible in the reported groups. The remaining terms describe distinctions inside groups, which that report no longer resolves. In the four-label example, all discrepancy lies in these conditional tables.

A whole family can collapse under grouping. For \[p_\theta=(\theta/2,(1-\theta)/2,\theta/2,(1-\theta)/2),
\qquad 0\leq\theta\leq1,\] the same grouping always outputs \((1/2,1/2)\). The parameter changes the fine distribution, but this observation scheme cannot identify it. This is a concrete information-geometric question: which directions through a family of probability models remain observable?

The conditional decomposition provides an equality condition for deterministic grouping. If the two models have identical conditional distributions inside every group carrying positive first-model mass, then grouping loses no KL divergence. Their entire distinction was already visible in the group totals.

For example, \(p=(0.3,0.3,0.2,0.2)\) and \(q=(0.1,0.1,0.4,0.4)\) differ in the totals of the two pairs, but both are uniform within each pair. Grouping the pairs retains their full KL divergence. This differs from the earlier example, where the group totals agreed and every distinction was internal.

The comparison gives a concrete way to design observations. If a question concerns relative mass between groups, recording group labels may retain the relevant distinction. If it concerns how probability is divided inside a group, that same recording scheme is insufficient. A suitable observation depends on the family of models and the question, rather than on the number of recorded symbols alone.

For the parameter family \(p_\theta\), grouping \(\{1,3\}\) and \(\{2,4\}\) returns \((\theta,1-\theta)\) and retains the parameter exactly. The alternative grouping returns a constant table. Both reports have two possible outputs; the decisive difference is which fine distinctions their outputs encode.

\section{Building a complete investigation}\label{ch10-building-a-complete-investigation}

A useful final investigation begins with a small transport graph and a precise update. Choose an initial probability vector and a step size satisfying the nonnegativity condition. Record several states, their differences, and their edge currents. Check conservation directly.

Next track entropy and divergence to a stationary reference. Separate a numerical pattern from a proof: the continuity equation proves conservation, double stochasticity proves entropy monotonicity, and a particular contraction estimate or exact solution can prove convergence.

Finally pass two initial models through the same observation channel. Compare their fine and reported divergences, overlaps, and spherical distances. If labels are grouped, use the conditional decomposition to identify the distinctions that disappear.

The \href{https://bobovski66.github.io/IntroDiscreteInfoGeom/demos/week10.html}{network demonstration} explores channels and model comparisons. The transport view of the \href{https://bobovski66.github.io/IntroDiscreteInfoGeom/demos/operator-explorer.html}{operator explorer} exposes currents and divergences directly. The \href{https://bobovski66.github.io/IntroDiscreteInfoGeom/guide.html\#capstone}{course guide} suggests longer investigations.

These calculations bring the course's original question into a concrete setting: a probability distribution changes through a sequence of finite operations, and the available observations determine which changes can still be distinguished.

\section{Exercises}\label{ch10-exercises}

\begin{exercises}

\item 

Derive the three-node Laplacian from its edge weights. Verify that \(I-\tfrac14L\) is the matrix used in Chapter 4.

\item 

Find a sufficient step-size bound for nonnegative diffusion on a four-vertex star with unit edge weights. Explain the role of its central degree.

\item 

For the three-node path and \(p=(1,0,0)\), calculate currents, divergence, and one probability update at \(\tau=1/4\).

\item 

Verify that \((1,0,-1)\) and \((1,-2,1)\) are multiplied by \(3/4\) and \(1/4\) under the path matrix.

\item 

Use the exact trajectory formula to calculate \(p_2\) from \(p_0=(1,0,0)\).

\item 

Prove entropy monotonicity for a doubly stochastic matrix using concavity of \(-x\ln x\).

\item 

Give a column-stochastic channel that lowers entropy. Explain why it does not contradict Exercise 6.

\item 

Derive the log-sum inequality from Gibbs' inequality by normalizing the two lists.

\item 

Prove data processing for \(D(Kp\|Kq)\), identifying where the channel column sums enter.

\item 

For the four-label grouping example, calculate input and output overlaps and compare Fisher--Rao distances.

\item 

Derive the conditional divergence decomposition for grouping. Use it to locate the lost distinction in the worked example.

\item 

Choose two distinct parameter values in the family \(p_\theta\). Show that the fine distributions differ but the grouped observations agree. Propose a different grouping that would retain information about \(\theta\).

\end{exercises}

\section{Selected solutions}\label{ch10-selected-solutions}

\textbf{2.} The central weighted degree is three and each outer degree is one. The bound is \(0\leq\tau\leq1/3\).

\textbf{3.} Currents are \(J_{AB}=1/4\) and \(J_{BC}=0\). Divergence is \((1/4,-1/4,0)\), giving next state \((3/4,1/4,0)\).

\textbf{5.} With \(a=1/2\), \(b=1/6\), the formula gives \((5/8,5/16,1/16)\), agreeing with direct multiplication.

\textbf{7.} The channel \(\left(\begin{smallmatrix}1&1\\0&0\end{smallmatrix}\right)\) sends a fair input to \((1,0)\), lowering entropy from \(\ln2\) to zero. Its first row sums to two, so it is not doubly stochastic.

\textbf{10.} Each input overlap term is \(\sqrt{0.04}=0.2\), so \(B(p,q)=0.8\) and \(d_{\mathrm{FR}}=2\arccos0.8\). Output overlap is one and output distance is zero.

\textbf{12.} Grouping \(\{1,3\}\) and \(\{2,4\}\) yields \((\theta,1-\theta)\), which retains the parameter. The original grouping always yields \((1/2,1/2)\).
\labheading\begin{enumerate}
\def\labelenumi{\arabic{enumi}.}
\tightlist
\item
  On A--B--C with \(p=(1,0,0)\) and \(\tau=1/4\), (a) compute both signed edge currents \(J_{AB},J_{BC}\), (b) calculate outward divergence at A, B, C, (c) use \(\Delta p=-\operatorname{div}J\) to determine the next state, and (d) check mass conservation.
\item
  Using the same path, calculate \(p_1\) through \(p_4\). (a) Write each state as a probability vector. (b) Compute all coordinate differences. (c) Verify every difference vector sums to zero. (d) Explain the approach toward uniform.
\item
  For \(p_0=(1,0,0)\) and the first three transport updates, (a) compute entropy at each step, (b) calculate its first differences, (c) identify the sign of every difference, and (d) explain which property of the diffusion matrix supports the pattern.
\item
  Compare two channels: the symmetric path diffusion and a channel sending every input to A. (a) Construct each stochastic matrix. (b) Calculate both output distributions from uniform input. (c) Compare input/output entropy. (d) Explain why a general stochastic channel does not have to raise entropy.
\item
  Capstone: choose a four- or five-node undirected graph and an initial probability distribution. (a) Define a legal diffusion step using \(\tau\leq1/d_{\max}\). (b) Calculate at least four iterates with a table or short program. (c) Show currents, divergence, mass conservation, and entropy changes. (d) Compare two states by Fisher--Rao distance and explain what the mathematical results say about the flow of information.
\end{enumerate}

\FloatBarrier

\end{document}
